\documentclass[../../../include/open-logic-section]{subfiles}

\begin{document}

\olfileid{sth}{spine}{valpha}

\olsection{将各个阶段定义为 $V_\alpha$}

在 \olref[sth][ordinals][]{chap} 中，我们定义了良序与 von Neumann（冯·诺伊曼）序数。在本章中，我们将用它们来刻画集合层级\emph{本身}。为此，回顾在 \olref[sth][ordinals][opps]{sec} 中，我们定义了后继序数与极限序数的概念。我们在以下定义中使用这些概念：

\begin{defn}\ollabel{defValphas}
	\begin{align*}
	V_\emptyset &\defis \emptyset\\
	V_{\ordsucc{\alpha}} &\defis \Pow{V_\alpha} & & 
	\text{对任意序数 }\alpha\\
	V_{\alpha} &\defis \bigcup_{\gamma < \alpha} V_\gamma & & 
	\text{当 }\alpha\text{ 是极限序数}
\end{align*}
\end{defn}

\noindent
这是通过序数上的\emph{超限递归}作出的定义。在这方面，我们应将其与自然数上函数的递归定义进行比较。\footnote{参见 \olref[sfr][infinite][dedekind]{sec} 中加法、乘法和指数的定义。}如同处理自然数时那样，需要定义基始情形与后继情形；但在处理序数时，我们还需要描述\emph{极限}情形的行为。

$V_\alpha$ 的这一定义将是一个重要的里程碑。我们曾非形式地指出，我们的集合层级是通过\emph{阶段}来形成集合的。$V_\alpha$ 实际上正是这些阶段。然而重要的是，这是对阶段的\emph{内在}刻画。我们并不是在提出该理论的一个可能\emph{模型}，而是在我们的集合论\emph{内部}定义这些阶段。

\end{document}